\documentclass[a4paper,11pt]{article}

% ---------------------------------------------------------------
% Encoding und Sprache (kompatibel mit xelatex)
% ---------------------------------------------------------------
\usepackage{fontspec}
\usepackage[ngerman,provide=*]{babel}
\usepackage{csquotes}

% ---------------------------------------------------------------
% Geometrie und Typographie
% ---------------------------------------------------------------
\usepackage[a4paper,margin=2.5cm,headheight=14pt]{geometry}
\usepackage{setspace}
\onehalfspacing
\usepackage{parskip}
\usepackage{booktabs}
\usepackage{tabularx}
\usepackage{array}
\usepackage{enumitem}
\setlist[itemize]{leftmargin=*,topsep=2pt,itemsep=2pt}
\setlist[enumerate]{leftmargin=*,topsep=2pt,itemsep=2pt}

% Kopf-/Fußzeile schlicht
\usepackage{fancyhdr}
\pagestyle{fancy}
\fancyhf{}
\renewcommand{\headrulewidth}{0.3pt}
\renewcommand{\footrulewidth}{0pt}
\fancyhead[L]{\footnotesize MTA -- Sitzung 1: Topic-Modelle als qualitative Methode}
\fancyhead[R]{\footnotesize Seite \thepage}
\fancyfoot[C]{\footnotesize Prof.\ Dr.\ Christian Papilloud · Institut für Soziologie · MLU Halle-Wittenberg}

% Schlichte Überschriften
\usepackage{titlesec}
\titleformat{\section}{\large\bfseries}{\thesection.}{0.6em}{}
\titleformat{\subsection}{\normalsize\bfseries}{\thesubsection}{0.6em}{}
\titleformat{\subsubsection}{\normalsize\itshape}{}{0pt}{}
\titlespacing*{\section}{0pt}{1.5em}{0.5em}
\titlespacing*{\subsection}{0pt}{1.0em}{0.3em}

% Hervorhebungen
\usepackage{xcolor}
\definecolor{boxgrau}{gray}{0.95}
\usepackage[most]{tcolorbox}
\newtcolorbox{merkbox}[1][]{
  colback=boxgrau,colframe=black!40,boxrule=0.4pt,
  arc=2pt,left=8pt,right=8pt,top=4pt,bottom=4pt,
  fonttitle=\bfseries,#1
}

\usepackage{hyperref}
\hypersetup{
  colorlinks=true,
  linkcolor=black!70,
  urlcolor=black!60,
  citecolor=black!60,
  pdftitle={MTA Sitzung 1 -- Handout},
  pdfauthor={Christian Papilloud},
}

% ---------------------------------------------------------------
% Titelinformationen
% ---------------------------------------------------------------
\title{\textbf{Topic-Modelle als qualitative Methode?}\\[4pt]
\large Sitzung 1 -- Positionierung gegenüber objektiver Hermeneutik
und qualitativer Inhaltsanalyse\\[8pt]
\normalsize Begleitendes Handout}
\author{Prof.\ Dr.\ Christian Papilloud\\
\small Institut für Soziologie -- Martin-Luther-Universität Halle-Wittenberg}
\date{Master Soziologie -- 2026}

\begin{document}

\maketitle
\thispagestyle{fancy}

\vspace{-1em}

% ===============================================================
\section{Worum es geht}

Bevor wir zur \emph{praktischen Anwendung} des Topic-Modell-Ansatzes
mit dem Programm \textbf{MTA} (Multi-Text Analyser) kommen,
beschäftigen wir uns in drei Sitzungen mit dem \emph{theoretischen
Hintergrund} von Topic-Modellen. Die heutige Sitzung beantwortet die
Grundfrage: \textbf{Ist Topic-Modellierung eine qualitative Methode
der Soziologie?} Wenn ja, in welchem Sinn?

Diese Frage ist nicht rein akademisch. Wer ein Topic-Modell-Programm
benutzt, ohne zu wissen, \emph{in welchen Methodentraditionen} dieses
Werkzeug steht, läuft Gefahr, seine Ergebnisse falsch zu
interpretieren -- entweder zu viel oder zu wenig in sie hineinzulesen.
Die Auseinandersetzung mit den klassischen qualitativen Methoden ist
deshalb keine Pflichtübung, sondern eine \emph{epistemologische
Vorbereitung} der Praxis.

% ===============================================================
\section{Die kontraintuitive These}

Topic-Modelle (TM) sind ein \emph{quantitativ basiertes} Werkzeug:
sie zählen Worthäufigkeiten, faktorisieren Matrizen und nutzen
Wahrscheinlichkeitsverteilungen. Sie dienen aber, das ist die These
des Buches, der \emph{qualitativen} Forschung. Diese These ist
kontraintuitiv, weil das Begriffspaar \enquote{quantitativ vs.\
qualitativ} in den Sozialwissenschaften meistens als unauflösbarer
Gegensatz auftritt: Statistik gegen Sinnverstehen, Zahlen gegen Texte,
Erklären gegen Verstehen.

\begin{merkbox}
\textbf{Worum es nicht geht.} Die These ist \emph{nicht}, dass
Topic-Modelle die Unterscheidung qualitativ/quantitativ überwinden
oder \emph{aufheben} -- diese Erwartung wird im Schlusskapitel des
Buches explizit zurückgewiesen (Kap.~8). Es geht vielmehr darum, zu
verstehen, dass ein quantitativ basiertes Verfahren \emph{qualitative
Erkenntnisaufgaben} erfüllen kann: Sinnstrukturen sichtbar machen,
die ein menschlicher Leser allein nicht erfassen könnte.
\end{merkbox}

% ===============================================================
\section{Drei gemeinsame Grundannahmen}

Bevor wir die Unterschiede betrachten, sollten wir die \emph{Gemein\-sam\-keiten}
ernst nehmen. Die objektive Hermeneutik, die qualitative
Inhaltsanalyse \emph{und} Topic-Modell-Verfahren gehören zum
Methodenarsenal der qualitativen Textanalyse, weil sie drei
Voraussetzungen teilen (Papilloud \& Hinneburg 2018, S.~2--3):

\begin{enumerate}
\item Begriffe und Sätze werden \textbf{nicht zufällig} verwendet.
Sinn ist kein zufälliges Ergebnis zusammengestellter Begriffe;
er wird durch ihren Gebrauch \emph{erzeugt}.

\item Texte enthalten \textbf{Strukturen}, die ihrem Sinn zugrunde
liegen. Die Aufgabe der Analyse besteht darin, diese Strukturen
zu rekonstruieren -- nicht, einen vorher feststehenden Inhalt
\enquote{wiederzugeben}.

\item Die \textbf{Deutung} -- und damit die Person, die deutet --
spielt eine konstitutive Rolle. Sowohl bei der Organisation des
analytischen Verfahrens als auch bei der Interpretation der Ergebnisse.
\end{enumerate}

Diese drei Annahmen reichen, um Topic-Modelle aus dem Lager der
\emph{rein deskriptiven} oder \emph{rein klassifikatorischen}
Statistik herauszuheben. TM sind nicht einfach \enquote{Wörter zählen}.

% ===============================================================
\section{Drei Methoden, drei Antworten}

Aus den gleichen Grundannahmen ergeben sich drei \emph{unterschiedliche}
Verfahren. Diese Sitzung kontrastiert sie. Sie unterscheiden sich an
drei Stellen:

\begin{itemize}
\item Was heißt \emph{Struktur} des Sinns? Wie ist sie aufgebaut?
\item Wie wird Sinn \emph{algorithmisch} erzeugt (im weiten Sinn von
\enquote{regelhaft})?
\item Welche \emph{Rolle} spielt der/die Forschende?
\end{itemize}

Wir gehen die drei Methoden nacheinander durch -- mit dem Schwerpunkt
auf das, was sie für den \emph{Kontrast} mit Topic-Modellen
interessant macht.

% ===============================================================
\section{Objektive Hermeneutik (Oevermann)}

Die objektive Hermeneutik wurde in den 1970er Jahren am Max-Planck-Institut
für Bildungsforschung entwickelt. Ulrich Oevermann, ihr wichtigster
Begründer, hat sie über mehrere Jahrzehnte ausgebaut. Sie ist
\emph{kein} reines Verfahren: sie ist mit einer spezifischen
\textbf{Theorie der Lebenspraxis} verbunden, in der jeder
Einzelakteur grundsätzlich \enquote{an jeder Sequenzstelle seines
Lebens hypothetische Welten} konstruiert (Oevermann 1996, 6).

\subsection{Lebenspraxis und objektive Sinnstruktur}

Lebenspraxis ist bei Oevermann das Zusammenwirken physischer,
psychischer, sozialer und kultureller Potentiale, die einem
Einzelakteur die Gestaltung seines Lebens erlauben (Oevermann 2004,
157). Das Leben wird als \textbf{Sequenz von Entscheidungen}
verstanden, in denen Sinnstrukturen \emph{gebildet} werden.

Ein zentraler Begriff ist die \textbf{Krise}: jede
Entscheidungssituation, in der der Akteur neue Handlungsregeln
gewinnt und erst \emph{retrospektiv} eine Begründung liefert. Es
gibt also eine \enquote{widersprüchliche Einheit von
Entscheidungszwang und Begründungsverpflichtung} (Oevermann
1996, 5--6): man muss sich entscheiden, ohne eine
Richtig-Falsch-Kalkulation zur Verfügung zu haben, und man muss
das Getane danach begründen.

Aus der Verkettung dieser Entscheidungen entsteht eine
\textbf{objektive Sinnstruktur}: eine \enquote{Fallstrukturgesetzlichkeit},
die der Lebenspraxis dieses bestimmten Individuums entspricht. Diese
Sinnstruktur soll die qualitative Analyse rekonstruieren.

\subsection{Manifeste und latente Sinnstrukturen}

Die objektive Sinnstruktur besteht aus zwei Schichten:

\begin{itemize}
\item Die \textbf{manifeste} Struktur: was ein Akteur explizit sagt,
was im Protokoll festgehalten ist -- empirisch beobachtbar.

\item Die \textbf{latente} Struktur: was sich aus der \emph{Verkettung}
der Aussagen ergibt. Sie ist nicht als konkrete Aussage protokolliert
und muss vom Forschenden \emph{gedeutet} werden (Oevermann 2002).
\end{itemize}

Hier liegt eine philosophisch starke Annahme: die latente Sinnstruktur
ist nicht das, was der Akteur \emph{wirklich gemeint} hat. Sie ist
nach Oevermann eine \enquote{abstrakte Realität}, die \enquote{von der
manifesten Realisierung dieser Sinnstrukturen im Bewusstsein eines
konkreten Autors oder Rezipienten unabhängig} ist (Oevermann 2002,
28--29). Diese Struktur wird durch \enquote{regelgeleitetes Handeln}
und durch \enquote{Regeln analog zu einem Algorithmus} erzeugt.

\subsubsection{Welcher Algorithmus?}

Der Begriff \enquote{Algorithmus} bei Oevermann ist \emph{nicht}
der mathematische Algorithmus eines Computers. Es ist die
generative Grammatik im Sinne Chomskys (1981): eine universelle
Menge von Regeln, die in der Sprache des Akteurs operieren und
Bedeutung erzeugen, ohne dass der Akteur sich dieser Regeln bewusst
ist. Wer spricht, wendet diese Regeln unbewusst an.

Diese Annahme ist die Quelle einer Reihe von Kritiken an der objektiven
Hermeneutik -- nicht zuletzt der Vorwurf, die Suche nach latenten
Sinnstrukturen sei eher ein \emph{metaphysisches} als ein empirisches
Problem (Reichertz 2004). Für uns ist nur wichtig festzuhalten:
\textbf{die objektive Hermeneutik setzt eine universelle generative
Grammatik voraus}. Wir werden sehen, dass dies eine zentrale Differenz
zu Topic-Modellen ist.

\subsection{Das Verfahren: Sequenzanalyse}

Konkret arbeitet die objektive Hermeneutik durch die
\textbf{Sequenzanalyse}. Das Protokoll wird Schritt für Schritt
durchgegangen. An jeder Sequenzstelle entwickelt der/die Forschende
viele \emph{Lesarten} -- Gedankenexperimente, alternative Deutungen,
die in diesem Kontext denkbar wären. Schrittweise werden Lesarten
eliminiert, die mit dem weiteren Verlauf des Protokolls nicht
vereinbar sind. Was übrig bleibt, ist die rekonstruierte
\textbf{Fallstrukturgesetzlichkeit} des beobachteten Falls.

Es ist ein arbeitsintensives Verfahren. Eine Sequenzanalyse kann
über Stunden hinweg \emph{einen einzigen Satz} betreffen. Sie wird
oft in Forschungsgruppen durchgeführt, damit verschiedene Lesarten
miteinander konfrontiert werden können. Typische Korpora sind
\emph{klein}: ein bis zwanzig Fälle.

% ===============================================================
\section{Qualitative Inhaltsanalyse (Mayring)}

Philipp Mayring entwickelt seit Anfang der 1980er Jahre einen
methodischen Vorschlag, der die qualitative Forschung von ihrer
\emph{handwerklichen} Variante in Richtung einer \emph{systematischen}
Auswertung bringen will. Sein Programm ist explizit \textbf{hybrid}:
es soll \enquote{die Vorteile der quantitativen Inhaltsanalyse bewahren
und auf qualitativ-interpretative Auswertungsschritte übertragen}.

\subsection{Das Kategoriensystem}

Die zentrale Operation ist die Konstruktion eines
\textbf{Kategoriensystems}. Es besteht aus:

\begin{itemize}
\item Kategorien und Unterkategorien;
\item Kategoriendefinitionen, die eindeutig sein müssen;
\item Ankerbeispielen, die zeigen, was zu einer Kategorie gehört;
\item einer Rückkopplungsschleife, in der die Kategorien im Laufe
der Auswertung überarbeitet werden.
\end{itemize}

Ein Kategoriensystem soll zwei Ansprüche miteinander verbinden:

\begin{enumerate}
\item In hermeneutischer Tradition: den \emph{manifesten} \emph{und}
den \emph{latenten} Sinn der Textquelle abbilden (Ramsenthaler 2013, 23).

\item In quantitativer Tradition: eine \emph{systematische Überprüfung}
von Strukturen ermöglichen, mit \textbf{Reliabilität},
\textbf{Inter\-ko\-der\-re\-lia\-bi\-li\-tät} und
\textbf{Triangulation} (Mayring 2000).
\end{enumerate}

Ziel ist eine \enquote{intersubjektiv nachvollziehbare} Auswertung
(Rinker 2013, 40): zwei unabhängige Codiererinnen sollen mit
demselben System zu vergleichbaren Ergebnissen kommen.

\subsection{Eine bekannte Kritik}

Mayrings Ansatz wird seit Jahrzehnten diskutiert und kritisiert.
Drei wiederkehrende Kritikpunkte sind:

\begin{enumerate}
\item Die \emph{Einbindung statistischer Methoden} bleibt
rudimentär. Sie beschränkt sich meistens auf das Auszählen von
Kategorienhäufigkeiten -- Kategorien, die zuvor ohnehin durch den
Forschenden codiert wurden.

\item Die Kategorien werden \emph{zwar} induktiv aus dem Text
gewonnen, sind \emph{aber} zunächst deduktiv angesetzt. Der Bezug
zum Text als Ganzes geht damit verloren (Plant 1996, Froggatt 2001,
Lamnek \& Krell 2005).

\item Nach Kruse (2015, 399) wirkt in manchen Spielarten eine
\emph{rein quantifizierende Subsumtionslogik}. Mayring würde
Äußerungen \emph{inventarisieren}, ohne die Ebene der
\emph{latenten Strukturen} zu erreichen. Damit blieben Fragen der
\textbf{Indexikalität} der Sprache unbeantwortet -- die Frage also,
wie Akteure ihre Perspektive(n) auf die Realität durch ihre
Sprachgebrauchspraxis konstruieren (Bohnsack 2003).
\end{enumerate}

Diese letzte Kritik ist für unsere Fragestellung besonders wichtig.
Sie wird uns helfen zu verstehen, was Topic-Modelle \emph{anders}
machen.

% ===============================================================
\section{Wo Topic-Modelle stehen}

Topic-Modelle teilen mit \emph{beiden} klassischen Methoden die drei
Grundannahmen vom Anfang. Sie unterscheiden sich aber in einem
entscheidenden Punkt:

\begin{merkbox}
\textbf{Was sich grundsätzlich ändert.} Bei Oevermann und Mayring ist
das, was die Analyse strukturiert -- die Fallstrukturgesetzlichkeit,
das Kategoriensystem --, das \emph{Ergebnis menschlicher Interpretation}.

Bei Topic-Modellen ist es das \emph{Ergebnis eines mathematischen
Algorithmus}, der dann seinerseits durch den Forschenden interpretiert
werden muss. Die Interpretation verschwindet nicht; sie verschiebt sich.
\end{merkbox}

\subsection{Drei Unterschiede zur objektiven Hermeneutik}

\paragraph{1. Begriff von \enquote{Algorithmus}.}
Bei Oevermann ist Algorithmus die regelgeleitete Selektion in
Möglichkeitswelten, die eine universelle generative Grammatik
voraussetzt. Bei Topic-Modellen sind Algorithmen
\emph{statistisch} oder \emph{linearealgebraisch} definiert. Sie
setzen \emph{nicht} voraus, dass die Grammatik universell ist; im
Gegenteil: die Regeln der Sprache verändern sich mit Gebrauch, Zeit,
Gesellschaft und Kultur (Foraker u.a.\ 2009, Kam \& Newport 2009,
Everett 2012).

\paragraph{2. Was wird gesucht?}
Bei Oevermann: die \emph{Fallstrukturgesetzlichkeit} eines bestimmten
Falls. Bei Topic-Modellen: \emph{Klassifikationsmuster} über viele
Dokumente hinweg. TM sind also strukturell auf den \textbf{Vergleich}
hin gebaut -- zwischen Texten, zwischen Akteuren, zwischen Zeitpunkten.

\paragraph{3. Rolle der Kontextvariation.}
Bei Oevermann: Gedankenexperimente und Lesarten werden \emph{aktiv}
mobilisiert; alternative Kontexte sind das Herz des Verfahrens. Bei
Topic-Modellen werden solche Variationen \emph{bewusst vermieden}
(Sinclair 1992). Es zählen die typischen Merkmale des Textes selbst.
Was interpretiert werden muss, ist das \textbf{Klassifikationsmuster},
das der Algorithmus erzeugt -- eine \emph{Hypothese} darüber, welche
Inhalte die Akteure mitgeteilt haben.

\subsection{Nähe zur qualitativen Inhaltsanalyse}

Topic-Modelle stehen Mayring deutlich näher als Oevermann. Beide
Verfahren:

\begin{itemize}
\item bauen Kategoriensysteme aus dem Material auf;
\item streben Intersubjektivität an;
\item nutzen Häufigkeiten als Grundlage.
\end{itemize}

Es gibt aber eine wichtige Differenz: was an Mayring kritisiert wird
(die Kategorien bleiben oberflächlich, sie inventarisieren, statt
latente Strukturen zu erreichen), versucht TM \emph{methodisch} zu
vermeiden. Topic-Modelle nehmen die \textbf{Indexikalität der Sprache}
ernst: Bedeutung ergibt sich aus der \emph{Konfiguration} von
Begriffen, nicht aus ihrer isolierten Codierung. Das, was die topics
ausmacht, ist die wiederholte Ko-Okkurrenz von Begriffen in
Dokumenten -- und das ist eine Aussage über die latente Struktur eines
Korpus, keine Inventarisierung der Oberfläche.

% ===============================================================
\section{Synopse}

Die folgende Tabelle fasst zusammen, was wir bisher gesehen haben.
Sie ist ein \emph{erstes} Schema; es wird in der dritten Sitzung
durch das, was Topic-Modelle nicht leisten können, präzisiert.

\smallskip

\begin{tabularx}{\textwidth}{@{}p{0.18\textwidth}XXX@{}}
\toprule
& \textbf{Objektive Hermeneutik} & \textbf{Qualitative Inhaltsanalyse} & \textbf{Topic-Modell} \\
\midrule
\textbf{Gegenstand}
  & Einzelfall, biographische Sequenz
  & Textmaterial, Inventar von Aussagen
  & Textsammlung, Korpus \\
\addlinespace[3pt]
\textbf{Algorithmus}
  & generative Grammatik (Chomsky)
  & Kategoriensystem, manuelle Codierung
  & statistische / lineare Algebra \\
\addlinespace[3pt]
\textbf{Sinnbegriff}
  & latent, regelerzeugt, objektiv
  & manifest \emph{und} latent
  & latent als Konfiguration von Begriffen \\
\addlinespace[3pt]
\textbf{Rolle der Forschenden}
  & Deutung, Entwicklung von Lesarten
  & Codierung anhand des Systems
  & Interpretation des Klassifikationsmusters \\
\addlinespace[3pt]
\textbf{Ergebnis}
  & Fallstrukturgesetzlichkeit
  & Kategorien mit Häufigkeiten
  & Topics + Dokument\-zugehörigkeiten \\
\addlinespace[3pt]
\textbf{Korpusgröße}
  & klein (1--20 Fälle)
  & mittel
  & groß (Hunderte bis Millionen) \\
\bottomrule
\end{tabularx}

% ===============================================================
\section{Fragen zum Mitnehmen}

Drei Fragen, über die Sie bis zur nächsten Sitzung nachdenken können:

\begin{enumerate}
\item Eine Sozialwissenschaftlerin, die mit MTA arbeitet, codiert nicht
selbst -- sie liest die Codierung einer Maschine. Verschiebt das die
\emph{Subjektivität der Interpretation} oder verlagert es sie nur an
einen anderen, weniger sichtbaren Ort?

\item Wir haben gesehen, dass Topic-Modelle dem Verfahren Mayrings
strukturell nahe stehen. Was, glauben Sie, würde Oevermann zu einer
Topic-Modell-Analyse einer Sammlung seiner eigenen Interviewprotokolle
sagen?

\item Die Behauptung, eine Methode könne \enquote{die latente
Sinnstruktur} eines Textes erschließen, wird sowohl von Oevermann als
auch (mit Einschränkungen) von Mayring erhoben. Topic-Modelle erheben
sie ebenfalls. Welche Form von \emph{Belastbarkeit} verlangen wir
jeweils, damit eine solche Behauptung legitim ist?
\end{enumerate}

% ===============================================================
\section{Ausblick auf Sitzung 2}

In der nächsten Sitzung verlassen wir die methodologische
Standortbestimmung und schauen uns an, \emph{wie Topic-Modelle
funktionieren}. Im Mittelpunkt:

\begin{itemize}
\item Was ist ein \emph{topic} mathematisch? Wir gehen das
Beispiel \enquote{MLU/Sommer} aus dem Buch durch (Abb.~2.1 in
Papilloud \& Hinneburg 2018).

\item Vergleich mit zwei verwandten Verfahren, die Sie aus
quantitativen Methoden kennen: Clusteranalyse und Faktorenanalyse.

\item Die zwei klassischen TM-Verfahren: \textbf{LDA} (Latent
Dirichlet Allocation, probabilistisch-bayessch) und \textbf{NMF}
(Non-negative Matrix Factorization, lineare Algebra). Beide
finden Sie als Optionen in MTA.

\item \textbf{Neue Welt:} \emph{BERTopic} (Grootendorst 2022) und
der Sprung zu \emph{Embeddings}. Was bringen Large Language Models
(GPT, Claude) in die Diskussion ein? Wir besprechen die jüngste
Literatur (2022--2026) zur Anwendung von LLMs in der qualitativen
Sozialforschung -- inklusive der \emph{Kritik} an dieser Praxis
(Ashwin u.a.\ 2025).

\item Eine Frage, die uns durch die Praxis mit MTA begleiten wird:
\emph{warum} verwenden wir NMF und LDA und nicht BERTopic oder
einen LLM? Es ist nicht nur eine technische Wahl.
\end{itemize}

% ===============================================================
\section*{Zitierte Literatur}
\addcontentsline{toc}{section}{Zitierte Literatur}

\footnotesize

\textsc{Grundtext} der Sitzung:

Papilloud, C., \& Hinneburg, A.\ (2018). \emph{Einführung in die
qualitative Analyse von Texten mit Topic-Modellen}. Wiesbaden:
Springer. \emph{Bes.\ Kap.~1.}

\bigskip

\textsc{Objektive Hermeneutik:}

Chomsky, N.\ (1981). \emph{Regeln und Repräsentationen}. Frankfurt:
Suhrkamp.

Oevermann, U.\ (1995). Ein Modell der Struktur von Religiosität.
In: M.\ Wohlrab-Sahr (Hrsg.), \emph{Biographie und Religion},
27--102. Frankfurt: Campus.

Oevermann, U.\ (1996). \emph{Krise und Muße. Struktureigenschaften
ästhetischer Erfahrung aus soziologischer Sicht}. Vortrag, Städelschule.

Oevermann, U.\ (2002). \emph{Klinische Soziologie auf der Basis der
Methodologie objektiver Hermeneutik -- Manifest}.

Oevermann, U.\ (2004). Sozialisation als Prozess der Krisenbewältigung.
In: D.\ Geulen \& H.\ Veith (Hrsg.), \emph{Sozialisationstheorie
interdisziplinär}, 155--181. Stuttgart: Lucius.

Oevermann, U.\ (2008). \emph{Krise und Routine als analytisches
Paradigma in den Sozialwissenschaften} (Abschiedsvorlesung).

Przyborski, A., \& Wohlrab-Sahr, M.\ (2009). \emph{Qualitative
Sozialforschung. Ein Arbeitsbuch}. München: Oldenbourg.

Reichertz, J.\ (2004). \enquote{Das Handlungsrepertoire von
Gesellschaften erweitern.} \emph{Forum Qualitative Sozialforschung}
5(3).

\bigskip

\textsc{Qualitative Inhaltsanalyse:}

Bohnsack, R.\ (2003). Dokumentarische Methode. In: R.\ Bohnsack
et al.\ (Hrsg.), \emph{Hauptbegriffe Qualitativer Sozialforschung},
40--44.

Kruse, J.\ (2015). \emph{Qualitative Interviewforschung. Ein
integrativer Ansatz}. Weinheim: Beltz.

Mayring, P.\ (1994). Qualitative Inhaltsanalyse. In: A.\ Boehm et al.\
(Hrsg.), \emph{Texte verstehen}, 159--175. Konstanz: UVK.

Mayring, P.\ (2000). Qualitative Inhaltsanalyse. \emph{Forum
Qualitative Sozialforschung} 1(2).

Mayring, P.\ (2002). \emph{Einführung in die qualitative
Sozialforschung}. Basel: Beltz.

Ramsenthaler, M.\ (2013). Was ist qualitative Inhaltsanalyse?
In: M.\ W.\ Schnell et al.\ (Hrsg.), \emph{Der Patient am Lebensende},
23--42. Wiesbaden: Springer.

Rinker, T.\ (2013). \emph{qdap: Quantitative Discourse Analysis
Package}. Buffalo: University at Buffalo.

\bigskip

\textsc{Topic-Modelle (Vorgriff auf Sitzung 2):}

Blei, D.\ M.\ (2012). Probabilistic Topic Models.
\emph{Communications of the ACM} 55(4): 77--84.

Grootendorst, M.\ (2022). BERTopic: Neural topic modeling with a
class-based TF-IDF procedure. \emph{arXiv:2203.05794}.

Sinclair, J.\ (1992). The Automatic Analysis of Corpora. In:
J.\ Svartvik (Hrsg.), \emph{Directions in Corpus Linguistics}.

\end{document}
